\section{Introduction}\label{sec:intro}

International product standards and recognition rules affect the cost of serving foreign markets and can determine whether a foreign producer remains an effective competitor. \citet{gandalshy2001} study these issues in a three-country model in which countries can recognize foreign standards or form a standardization union. Within a union, member firms sell without conversion costs while an outsider faces a conversion cost in member markets. The model provides a compact benchmark for standards institutions in international trade; related work studies alternative recognition and compatibility arrangements \citep{costinot2008,klimenko2009}. This note revisits the price subgame after the outsider is foreclosed in a union-member market and asks a narrower question than a new standards-policy model: what equilibrium correspondence is implied by the published quadratic-transport specification, and what does that imply for the original welfare comparison?

The foreclosure threshold itself is not in dispute. In the published three-firm interior solution, outsider demand reaches zero at $c=5/2$. The problem arises after foreclosure. The two surviving member firms then compete across both the short and long arcs joining their locations. Under squared-distance transportation costs, the long-arc marginal consumer responds to a price difference with coefficient $1/4$, whereas Appendix B uses the same $1/2$ coefficient as on the short arc. The resulting complete published profile $(3/2,3/2,c)$ is not Nash in the strict post-foreclosure range: a sufficiently small member price increase is profitable. The correction therefore concerns the published price profile and its uniqueness claim, not the foreclosure threshold and not the proposition that a member price of $3/2$ can never arise under another outsider quote.

Correcting the geometry reveals equilibrium multiplicity rather than a single replacement price. Within symmetric, foreclosed, pure-strategy profiles in one member market, the unrestricted original price game has a continuum of equilibria. For member price $s<2$, a zero-sales outsider quote at $s+1$ changes the members' global upward-deviation incentives and can support any $s$ between $3/2$ and the relevant cost bound. At $s=2$, the usual two-member optimum is also supportable when conversion cost is high enough. The multiplicity is economically relevant because the outsider's posted price can discipline active firms even when its realized sales are zero.

A separate benchmark clarifies the piecewise price formula obtained in the earlier correction attempt. If the strategy set is explicitly restricted to prices no lower than marginal cost in each market, then within the same symmetric, foreclosed, pure-strategy class the member price is pinned down at $c-1$ for $5/2<c<3$ and at $2$ for $3\leq c<5$. This no-below-cost restriction is a modification of the strategy space; it is not stated in \citet{gandalshy2001}, is not without loss of generality, and is not derived here from weak dominance. On its lower branch the outsider has zero sales but its cost-based exclusion constraint binds member pricing, a static exclusion-maintenance mechanism that can be described as limit pricing in the product-market sense rather than the signaling sense of \citet{milgromroberts1982}.

The revised equilibrium correspondence also changes the scope of the welfare statement. If both segmented union markets select the same symmetric continuation price, the member-country price terms cancel: national welfare remains $3V+1/4$, one half above the mutual-recognition benchmark used in Proposition 3 of \citet{gandalshy2001}. If the two union markets select different symmetric continuation prices, however, the domestic firm's export-price term and domestic consumer-price term no longer cancel. The original ranking is therefore preserved under a common symmetric continuation, including the symmetric continuation in the no-below-cost benchmark, but it is not selection-free across arbitrary unrestricted market-by-market continuations.

The contribution is a correction-plus-equilibrium-selection result. It establishes that the published post-foreclosure profile is not Nash, characterizes the unrestricted symmetric foreclosed pure-strategy equilibrium set, shows exactly what an explicit no-below-cost restriction does within that class, and states the welfare comparison at the corresponding continuation-selection scope. It does not characterize all asymmetric or mixed equilibria, solve the government stage under unrestricted multiplicity, claim that the original policy ranking reverses, or claim general novelty for foreclosure or limit pricing. The purpose is instead to delimit precisely what the published benchmark implies once its post-foreclosure demand geometry is made consistent with its quadratic transportation primitive.
